% --- Archon physics formalization source begin ---
% archon:physics
% archon:covers PhyXMiniProblems/problem_phyx_mini_0692.lean
% archon:source-report reports/phyx_mini/problem_phyx_mini_0692.source.json

\chapter{Physics problem phyx\_mini\_0692}
\label{ch:PhyXMiniProblems_problem_phyx_mini_0692}

\paragraph{Problem source.}
\#\# Physical scenario

A pendulum with a cord of length \textbackslash{}( r = 1.00 \textbackslash{}text\{ m\} \textbackslash{}) swings in a vertical plane as shown in figure. When the pendulum is in the two horizontal positions \textbackslash{}( \textbackslash{}theta = 90.0\textasciicircum{}\{\textbackslash{}circ\} \textbackslash{}) and \textbackslash{}( \textbackslash{}theta = 270\textasciicircum{}\{\textbackslash{}circ\} \textbackslash{}), its speed is \textbackslash{}( 5.00 \textbackslash{}text\{ m/s\} \textbackslash{}).

\#\# Question

Calculate the magnitude of the total acceleration at these positions.

\#\# Answer choices

- A: A: \textbackslash{}( 14.5 \textbackslash{}text\{ m/s\}\textasciicircum{}2 \textbackslash{})

- B: B: \textbackslash{}( 12.8 \textbackslash{}text\{ m/s\}\textasciicircum{}2 \textbackslash{})

- C: C: \textbackslash{}( 26.8 \textbackslash{}text\{ m/s\}\textasciicircum{}2 \textbackslash{})

- D: D: \textbackslash{}( 21.1 \textbackslash{}text\{ m/s\}\textasciicircum{}2 \textbackslash{})

\#\# Figure caption (auxiliary; use the image as primary evidence)

The image depicts a pendulum system with several components and vector notations:

- A mass (represented by a ball) is suspended by a string from a fixed point.
- The string makes an angle \textbackslash{}(\textbackslash{}theta\textbackslash{}) with the vertical dashed line.
- The pendulum's path is outlined by a dashed arc, showing it swings in a circular motion.
- There are three vectors emanating from the ball:
  - \textbackslash{}(\textbackslash{}vec\{a\}\textbackslash{}), representing the total acceleration, points diagonally upwards.
  - \textbackslash{}(\textbackslash{}vec\{a\}\_r\textbackslash{}), the radial acceleration, points along the string towards the pivot.
  - \textbackslash{}(\textbackslash{}vec\{a\}\_t\textbackslash{}), the tangential acceleration, points perpendicular to the radial direction along the path.
- The angle \textbackslash{}(\textbackslash{}phi\textbackslash{}) is formed between \textbackslash{}(\textbackslash{}vec\{a\}\textbackslash{}) and \textbackslash{}(\textbackslash{}vec\{a\}\_r\textbackslash{}).
- The downward vector \textbackslash{}(\textbackslash{}vec\{g\}\textbackslash{}) represents gravitational acceleration.
- The string length is denoted by \textbackslash{}(r\textbackslash{}).

Overall, the diagram illustrates the forces and accelerations acting on a pendulum in motion.

\#\# Dataset metadata

Rotational Motion; Physical Model Grounding Reasoning, Multi-Formula Reasoning

\paragraph{Recorded answer/context.}
C: C: \textbackslash{}( 26.8 \textbackslash{}text\{ m/s\}\textasciicircum{}2 \textbackslash{})

\paragraph{Figure/image path.}
/root/proposal\_for\_physic/science-mango/phyx\_mini\_run/phyx\_data/test\_image/692.png

\paragraph{Formalization target.}
create a compiling Lean file with sorry bodies at `PhyXMiniProblems/problem\_phyx\_mini\_0692.lean`. The Lean declarations must preserve the physical quantities, dimensions or dimensional roles, figure labels, governing-law hypotheses, and final relation expressed by this problem.
Use Mathlib/Physlib names found through LeanExplore where available. If a physics API is missing, introduce faithful local abstractions rather than scalar placeholder aliases.

\begin{theorem}[Physics formalization target]
\label{thm:physics:phyx_mini_0692:target}
\lean{PhyXMiniProblems.ProblemPhyXMini0692.totalAccelerationAtHorizontalPositions}
\uses{def:physics:phyx-mini-0692:phyxminiproblems-problemphyxmini0692-accelerationmagnitudeinmeterspersecondsquared, def:physics:phyx-mini-0692:phyxminiproblems-problemphyxmini0692-horizontalposition, def:physics:phyx-mini-0692:phyxminiproblems-problemphyxmini0692-pendulumsetup, def:physics:phyx-mini-0692:phyxminiproblems-problemphyxmini0692-matchesproblemstatement, def:physics:phyx-mini-0692:phyxminiproblems-problemphyxmini0692-usesstandardterrestrialgravity, def:physics:phyx-mini-0692:phyxminiproblems-problemphyxmini0692-hasphysicalpendulumparameters, def:physics:phyx-mini-0692:phyxminiproblems-problemphyxmini0692-matchessuppliedfigure, def:physics:phyx-mini-0692:phyxminiproblems-problemphyxmini0692-satisfieshorizontalpendulumlaws, def:physics:phyx-mini-0692:phyxminiproblems-problemphyxmini0692-isclosestdisplayedanswer}
The assigned autoformalize agent should translate this physics problem into Lean declarations in the covered file, with theorem and lemma proof bodies written as `by sorry`.
\end{theorem}
\begin{proof}
This is an autoformalization task, not a proof task. Produce statements that can later be proved by the physics prover without weakening the physical model.
\end{proof}

% --- Archon declaration topology begin ---
\section{Declaration topology}
\begin{definition}[acceleration Dimension]
  \label{def:physics:phyx-mini-0692:phyxminiproblems-problemphyxmini0692-accelerationdimension}
  \lean{PhyXMiniProblems.ProblemPhyXMini0692.accelerationDimension}
  The physical dimension L T⁻² of acceleration
\end{definition}
\begin{proof}
  This is the declared model definition of acceleration Dimension.
\end{proof}
\begin{definition}[Pendulum Plane]
  \label{def:physics:phyx-mini-0692:phyxminiproblems-problemphyxmini0692-pendulumplane}
  \lean{PhyXMiniProblems.ProblemPhyXMini0692.PendulumPlane}
  The two-dimensional vertical plane in which the pendulum swings
\end{definition}
\begin{proof}
  This is the declared model definition of Pendulum Plane.
\end{proof}
\begin{definition}[Length Quantity]
  \label{def:physics:phyx-mini-0692:phyxminiproblems-problemphyxmini0692-lengthquantity}
  \lean{PhyXMiniProblems.ProblemPhyXMini0692.LengthQuantity}
  A unit-independent physical length
\end{definition}
\begin{proof}
  This is the declared model definition of Length Quantity.
\end{proof}
\begin{definition}[Speed Quantity]
  \label{def:physics:phyx-mini-0692:phyxminiproblems-problemphyxmini0692-speedquantity}
  \lean{PhyXMiniProblems.ProblemPhyXMini0692.SpeedQuantity}
  A unit-independent physical speed
\end{definition}
\begin{proof}
  This is the declared model definition of Speed Quantity.
\end{proof}
\begin{definition}[Planar Position Quantity]
  \label{def:physics:phyx-mini-0692:phyxminiproblems-problemphyxmini0692-planarpositionquantity}
  \lean{PhyXMiniProblems.ProblemPhyXMini0692.PlanarPositionQuantity}
  \uses{def:physics:phyx-mini-0692:phyxminiproblems-problemphyxmini0692-pendulumplane}
  A unit-independent point or displacement vector in the pendulum plane
\end{definition}
\begin{proof}
  This is the declared model definition of Planar Position Quantity.
\end{proof}
\begin{definition}[Planar Acceleration Quantity]
  \label{def:physics:phyx-mini-0692:phyxminiproblems-problemphyxmini0692-planaraccelerationquantity}
  \lean{PhyXMiniProblems.ProblemPhyXMini0692.PlanarAccelerationQuantity}
  \uses{def:physics:phyx-mini-0692:phyxminiproblems-problemphyxmini0692-accelerationdimension, def:physics:phyx-mini-0692:phyxminiproblems-problemphyxmini0692-pendulumplane}
  A unit-independent acceleration vector in the pendulum plane
\end{definition}
\begin{proof}
  This is the declared model definition of Planar Acceleration Quantity.
\end{proof}
\begin{definition}[length Readout]
  \label{def:physics:phyx-mini-0692:phyxminiproblems-problemphyxmini0692-lengthreadout}
  \lean{PhyXMiniProblems.ProblemPhyXMini0692.lengthReadout}
  \uses{def:physics:phyx-mini-0692:phyxminiproblems-problemphyxmini0692-lengthquantity}
  Scalar readout of a physical length in a coherent choice of units
\end{definition}
\begin{proof}
  This is the declared model definition of length Readout.
\end{proof}
\begin{definition}[speed Readout]
  \label{def:physics:phyx-mini-0692:phyxminiproblems-problemphyxmini0692-speedreadout}
  \lean{PhyXMiniProblems.ProblemPhyXMini0692.speedReadout}
  \uses{def:physics:phyx-mini-0692:phyxminiproblems-problemphyxmini0692-speedquantity}
  Scalar readout of a physical speed in a coherent choice of units
\end{definition}
\begin{proof}
  This is the declared model definition of speed Readout.
\end{proof}
\begin{definition}[position Vector Readout]
  \label{def:physics:phyx-mini-0692:phyxminiproblems-problemphyxmini0692-positionvectorreadout}
  \lean{PhyXMiniProblems.ProblemPhyXMini0692.positionVectorReadout}
  \uses{def:physics:phyx-mini-0692:phyxminiproblems-problemphyxmini0692-pendulumplane, def:physics:phyx-mini-0692:phyxminiproblems-problemphyxmini0692-planarpositionquantity}
  Planar readout of a physical position in a coherent choice of units
\end{definition}
\begin{proof}
  This is the declared model definition of position Vector Readout.
\end{proof}
\begin{definition}[acceleration Vector Readout]
  \label{def:physics:phyx-mini-0692:phyxminiproblems-problemphyxmini0692-accelerationvectorreadout}
  \lean{PhyXMiniProblems.ProblemPhyXMini0692.accelerationVectorReadout}
  \uses{def:physics:phyx-mini-0692:phyxminiproblems-problemphyxmini0692-pendulumplane, def:physics:phyx-mini-0692:phyxminiproblems-problemphyxmini0692-planaraccelerationquantity}
  Planar readout of an acceleration vector in coherent acceleration units
\end{definition}
\begin{proof}
  This is the declared model definition of acceleration Vector Readout.
\end{proof}
\begin{definition}[acceleration Magnitude Readout]
  \label{def:physics:phyx-mini-0692:phyxminiproblems-problemphyxmini0692-accelerationmagnitudereadout}
  \lean{PhyXMiniProblems.ProblemPhyXMini0692.accelerationMagnitudeReadout}
  \uses{def:physics:phyx-mini-0692:phyxminiproblems-problemphyxmini0692-planaraccelerationquantity, def:physics:phyx-mini-0692:phyxminiproblems-problemphyxmini0692-accelerationvectorreadout}
  Magnitude readout of a physical acceleration vector
\end{definition}
\begin{proof}
  This is the declared model definition of acceleration Magnitude Readout.
\end{proof}
\begin{definition}[length In Meters]
  \label{def:physics:phyx-mini-0692:phyxminiproblems-problemphyxmini0692-lengthinmeters}
  \lean{PhyXMiniProblems.ProblemPhyXMini0692.lengthInMeters}
  \uses{def:physics:phyx-mini-0692:phyxminiproblems-problemphyxmini0692-lengthquantity, def:physics:phyx-mini-0692:phyxminiproblems-problemphyxmini0692-lengthreadout}
  Length readout in SI metres
\end{definition}
\begin{proof}
  This is the declared model definition of length In Meters.
\end{proof}
\begin{definition}[speed In Meters Per Second]
  \label{def:physics:phyx-mini-0692:phyxminiproblems-problemphyxmini0692-speedinmeterspersecond}
  \lean{PhyXMiniProblems.ProblemPhyXMini0692.speedInMetersPerSecond}
  \uses{def:physics:phyx-mini-0692:phyxminiproblems-problemphyxmini0692-speedquantity, def:physics:phyx-mini-0692:phyxminiproblems-problemphyxmini0692-speedreadout}
  Speed readout in SI metres per second
\end{definition}
\begin{proof}
  This is the declared model definition of speed In Meters Per Second.
\end{proof}
\begin{definition}[acceleration Magnitude In Meters Per Second Squared]
  \label{def:physics:phyx-mini-0692:phyxminiproblems-problemphyxmini0692-accelerationmagnitudeinmeterspersecondsquared}
  \lean{PhyXMiniProblems.ProblemPhyXMini0692.accelerationMagnitudeInMetersPerSecondSquared}
  \uses{def:physics:phyx-mini-0692:phyxminiproblems-problemphyxmini0692-planaraccelerationquantity, def:physics:phyx-mini-0692:phyxminiproblems-problemphyxmini0692-accelerationmagnitudereadout}
  Acceleration-magnitude readout in SI metres per second squared
\end{definition}
\begin{proof}
  This is the declared model definition of acceleration Magnitude In Meters Per Second Squared.
\end{proof}
\begin{definition}[Horizontal Position]
  \label{def:physics:phyx-mini-0692:phyxminiproblems-problemphyxmini0692-horizontalposition}
  \lean{PhyXMiniProblems.ProblemPhyXMini0692.HorizontalPosition}
  The two horizontal positions named in the problem
\end{definition}
\begin{proof}
  This is the declared model definition of Horizontal Position.
\end{proof}
\begin{definition}[theta Degrees]
  \label{def:physics:phyx-mini-0692:phyxminiproblems-problemphyxmini0692-horizontalposition-thetadegrees}
  \lean{PhyXMiniProblems.ProblemPhyXMini0692.HorizontalPosition.thetaDegrees}
  \uses{def:physics:phyx-mini-0692:phyxminiproblems-problemphyxmini0692-horizontalposition}
  The directed degree label attached to each horizontal position
\end{definition}
\begin{proof}
  This is the declared model definition of theta Degrees.
\end{proof}
\begin{definition}[Figure Element]
  \label{def:physics:phyx-mini-0692:phyxminiproblems-problemphyxmini0692-figureelement}
  \lean{PhyXMiniProblems.ProblemPhyXMini0692.FigureElement}
  Geometric objects and vector labels visible in the supplied bitmap
\end{definition}
\begin{proof}
  This is the declared model definition of Figure Element.
\end{proof}
\begin{definition}[Pendulum Figure]
  \label{def:physics:phyx-mini-0692:phyxminiproblems-problemphyxmini0692-pendulumfigure}
  \lean{PhyXMiniProblems.ProblemPhyXMini0692.PendulumFigure}
  \uses{def:physics:phyx-mini-0692:phyxminiproblems-problemphyxmini0692-figureelement}
  Presentation-level visibility data transcribed from the primary image
\end{definition}
\begin{proof}
  This is the declared model definition of Pendulum Figure.
\end{proof}
\begin{definition}[Horizontal Pendulum State]
  \label{def:physics:phyx-mini-0692:phyxminiproblems-problemphyxmini0692-horizontalpendulumstate}
  \lean{PhyXMiniProblems.ProblemPhyXMini0692.HorizontalPendulumState}
  \uses{def:physics:phyx-mini-0692:phyxminiproblems-problemphyxmini0692-speedquantity, def:physics:phyx-mini-0692:phyxminiproblems-problemphyxmini0692-planarpositionquantity, def:physics:phyx-mini-0692:phyxminiproblems-problemphyxmini0692-planaraccelerationquantity}
  The kinematic state at one horizontal position. The fields correspond to the image labels theta, a, a\_r, a\_t, g, and phi. In particular, the total acceleration remains an independent physical vector; no numerical answer is stored in this structure
\end{definition}
\begin{proof}
  This is the declared model definition of Horizontal Pendulum State.
\end{proof}
\begin{definition}[Pendulum Setup]
  \label{def:physics:phyx-mini-0692:phyxminiproblems-problemphyxmini0692-pendulumsetup}
  \lean{PhyXMiniProblems.ProblemPhyXMini0692.PendulumSetup}
  \uses{def:physics:phyx-mini-0692:phyxminiproblems-problemphyxmini0692-pendulumplane, def:physics:phyx-mini-0692:phyxminiproblems-problemphyxmini0692-lengthquantity, def:physics:phyx-mini-0692:phyxminiproblems-problemphyxmini0692-planarpositionquantity, def:physics:phyx-mini-0692:phyxminiproblems-problemphyxmini0692-horizontalposition, def:physics:phyx-mini-0692:phyxminiproblems-problemphyxmini0692-pendulumfigure, def:physics:phyx-mini-0692:phyxminiproblems-problemphyxmini0692-horizontalpendulumstate}
  The fixed geometry, length, two horizontal states, and supplied figure
\end{definition}
\begin{proof}
  This is the declared model definition of Pendulum Setup.
\end{proof}
\begin{definition}[Matches Problem Statement]
  \label{def:physics:phyx-mini-0692:phyxminiproblems-problemphyxmini0692-matchesproblemstatement}
  \lean{PhyXMiniProblems.ProblemPhyXMini0692.MatchesProblemStatement}
  \uses{def:physics:phyx-mini-0692:phyxminiproblems-problemphyxmini0692-lengthinmeters, def:physics:phyx-mini-0692:phyxminiproblems-problemphyxmini0692-speedinmeterspersecond, def:physics:phyx-mini-0692:phyxminiproblems-problemphyxmini0692-pendulumsetup}
  The numerical data stated in the prose: r = 1.00 m, both horizontal speeds are 5.00 m/s, and the two directed position labels are 90 and 270 degrees. No total-acceleration value occurs here
\end{definition}
\begin{proof}
  This is the declared model definition of Matches Problem Statement.
\end{proof}
\begin{definition}[Uses Standard Terrestrial Gravity]
  \label{def:physics:phyx-mini-0692:phyxminiproblems-problemphyxmini0692-usesstandardterrestrialgravity}
  \lean{PhyXMiniProblems.ProblemPhyXMini0692.UsesStandardTerrestrialGravity}
  \uses{def:physics:phyx-mini-0692:phyxminiproblems-problemphyxmini0692-accelerationmagnitudeinmeterspersecondsquared, def:physics:phyx-mini-0692:phyxminiproblems-problemphyxmini0692-pendulumsetup}
  The standard terrestrial-gravity calibration implicit in the numerical multiple-choice problem. This supplies 9.80 m/s², not the requested total acceleration
\end{definition}
\begin{proof}
  This is the declared model definition of Uses Standard Terrestrial Gravity.
\end{proof}
\begin{definition}[Has Physical Pendulum Parameters]
  \label{def:physics:phyx-mini-0692:phyxminiproblems-problemphyxmini0692-hasphysicalpendulumparameters}
  \lean{PhyXMiniProblems.ProblemPhyXMini0692.HasPhysicalPendulumParameters}
  \uses{def:physics:phyx-mini-0692:phyxminiproblems-problemphyxmini0692-lengthinmeters, def:physics:phyx-mini-0692:phyxminiproblems-problemphyxmini0692-speedinmeterspersecond, def:physics:phyx-mini-0692:phyxminiproblems-problemphyxmini0692-accelerationmagnitudeinmeterspersecondsquared, def:physics:phyx-mini-0692:phyxminiproblems-problemphyxmini0692-pendulumsetup}
  Positivity and normalization conditions for the nondegenerate setup
\end{definition}
\begin{proof}
  This is the declared model definition of Has Physical Pendulum Parameters.
\end{proof}
\begin{definition}[Matches Supplied Figure]
  \label{def:physics:phyx-mini-0692:phyxminiproblems-problemphyxmini0692-matchessuppliedfigure}
  \lean{PhyXMiniProblems.ProblemPhyXMini0692.MatchesSuppliedFigure}
  \uses{def:physics:phyx-mini-0692:phyxminiproblems-problemphyxmini0692-lengthreadout, def:physics:phyx-mini-0692:phyxminiproblems-problemphyxmini0692-positionvectorreadout, def:physics:phyx-mini-0692:phyxminiproblems-problemphyxmini0692-accelerationvectorreadout, def:physics:phyx-mini-0692:phyxminiproblems-problemphyxmini0692-pendulumsetup}
  Geometry read from the image: every named element is present, the bob is one cord length from the pivot, a\_r points toward the pivot, a\_r and a\_t are perpendicular, g points downward, and phi is the angle from a\_r to a. These are geometric readouts, not a numerical total-acceleration answer
\end{definition}
\begin{proof}
  This is the declared model definition of Matches Supplied Figure.
\end{proof}
\begin{definition}[Satisfies Horizontal Pendulum Laws]
  \label{def:physics:phyx-mini-0692:phyxminiproblems-problemphyxmini0692-satisfieshorizontalpendulumlaws}
  \lean{PhyXMiniProblems.ProblemPhyXMini0692.SatisfiesHorizontalPendulumLaws}
  \uses{def:physics:phyx-mini-0692:phyxminiproblems-problemphyxmini0692-lengthreadout, def:physics:phyx-mini-0692:phyxminiproblems-problemphyxmini0692-speedreadout, def:physics:phyx-mini-0692:phyxminiproblems-problemphyxmini0692-accelerationvectorreadout, def:physics:phyx-mini-0692:phyxminiproblems-problemphyxmini0692-accelerationmagnitudereadout, def:physics:phyx-mini-0692:phyxminiproblems-problemphyxmini0692-pendulumsetup}
  The physical laws used at either horizontal position: total acceleration is the vector sum a = a\_r + a\_t; circular-motion radial acceleration has magnitude v²/r; at a horizontal position gravity is wholly tangential. The orthogonality of the displayed radial and tangential arrows is retained separately in MatchesSuppliedFigure. None of these laws states the requested total magnitude or an answer choice
\end{definition}
\begin{proof}
  This is the declared model definition of Satisfies Horizontal Pendulum Laws.
\end{proof}
\begin{lemma}[total Acceleration Magnitude Formula]
  \label{lem:physics:phyx-mini-0692:phyxminiproblems-problemphyxmini0692-totalaccelerationmagnitudeformula}
  \lean{PhyXMiniProblems.ProblemPhyXMini0692.totalAccelerationMagnitudeFormula}
  \uses{def:physics:phyx-mini-0692:phyxminiproblems-problemphyxmini0692-lengthreadout, def:physics:phyx-mini-0692:phyxminiproblems-problemphyxmini0692-speedreadout, def:physics:phyx-mini-0692:phyxminiproblems-problemphyxmini0692-accelerationmagnitudereadout, def:physics:phyx-mini-0692:phyxminiproblems-problemphyxmini0692-horizontalposition, def:physics:phyx-mini-0692:phyxminiproblems-problemphyxmini0692-pendulumsetup, def:physics:phyx-mini-0692:phyxminiproblems-problemphyxmini0692-hasphysicalpendulumparameters, def:physics:phyx-mini-0692:phyxminiproblems-problemphyxmini0692-matchessuppliedfigure, def:physics:phyx-mini-0692:phyxminiproblems-problemphyxmini0692-satisfieshorizontalpendulumlaws}
  Pythagorean composition of the radial and tangential acceleration components. This generic lemma leaves the problem-specific numerical substitutions to the main theorem
\end{lemma}
\begin{proof}
  The conclusion follows from the governing relations and figure readouts named in the statement of total Acceleration Magnitude Formula.
\end{proof}
\begin{definition}[Answer Choice]
  \label{def:physics:phyx-mini-0692:phyxminiproblems-problemphyxmini0692-answerchoice}
  \lean{PhyXMiniProblems.ProblemPhyXMini0692.AnswerChoice}
  Labels of the four displayed multiple-choice answers
\end{definition}
\begin{proof}
  This is the declared model definition of Answer Choice.
\end{proof}
\begin{definition}[answer Magnitude In Meters Per Second Squared]
  \label{def:physics:phyx-mini-0692:phyxminiproblems-problemphyxmini0692-answermagnitudeinmeterspersecondsquared}
  \lean{PhyXMiniProblems.ProblemPhyXMini0692.answerMagnitudeInMetersPerSecondSquared}
  \uses{def:physics:phyx-mini-0692:phyxminiproblems-problemphyxmini0692-answerchoice}
  Displayed acceleration-magnitude values, in metres per second squared
\end{definition}
\begin{proof}
  This is the declared model definition of answer Magnitude In Meters Per Second Squared.
\end{proof}
\begin{definition}[Is Closest Displayed Answer]
  \label{def:physics:phyx-mini-0692:phyxminiproblems-problemphyxmini0692-isclosestdisplayedanswer}
  \lean{PhyXMiniProblems.ProblemPhyXMini0692.IsClosestDisplayedAnswer}
  \uses{def:physics:phyx-mini-0692:phyxminiproblems-problemphyxmini0692-answerchoice, def:physics:phyx-mini-0692:phyxminiproblems-problemphyxmini0692-answermagnitudeinmeterspersecondsquared}
  A choice is the unique displayed value closest to an exact SI magnitude
\end{definition}
\begin{proof}
  This is the declared model definition of Is Closest Displayed Answer.
\end{proof}
% --- Archon declaration topology end ---
% --- Archon physics formalization source end ---
